% Phongsakon Mark Konrad, academic research CV
\documentclass[10pt,a4paper,sans]{moderncv}
\moderncvstyle{classic}
\moderncvcolor{black}
\usepackage[T1]{fontenc}
\usepackage[utf8]{inputenc}
\usepackage[a4paper,top=1.5cm,bottom=1.5cm,left=1.5cm,right=1.5cm]{geometry}
\usepackage{enumitem}
\usepackage{lastpage}
\usepackage{needspace}
\setlength{\hintscolumnwidth}{2.15cm}
\setlength{\footskip}{12pt}
\setlength{\parskip}{0pt}
\setlist[itemize]{leftmargin=1.2em,itemsep=1pt,topsep=2pt,parsep=0pt}
% Small icon links for publications
\newcommand{\iconlink}[3]{\href{#2}{{\footnotesize #1}\,#3}}
\AtBeginDocument{\hypersetup{colorlinks=true,urlcolor=black,linkcolor=black,citecolor=black}}

\name{Phongsakon Mark}{Konrad}
\title{Curriculum Vitae}
\address{T\"ubingen, Germany}{}{}
\email{phongsakon@outlook.dk}
\homepage{phomarkon.github.io}
\social[linkedin]{phongsakonmarkkonrad}
\social[github]{phomarkon}
\social[orcid]{0009-0004-2521-7879}
\extrainfo{\protect\href{https://scholar.google.com/citations?user=zPHjLREAAAAJ}{Google Scholar}}

\begin{document}
\fancyfoot[l]{\color{color2}\footnotesize Last updated 23 November 2026}
\fancyfoot[r]{\smash{\raisebox{5pt}{\color{color2}\pagenumberfont\thepage/\protect\NoHyper\pageref{LastPage}\protect\endNoHyper}}}
\makecvtitle
\vspace{-7mm}

\section{Summary}
\cvitem{}{I work on interpretability. My question is whether we can trust what we see from a model (what it says about itself, how it behaves when tested, what we can read out of it), and I test this by changing things inside the model rather than by looking at correlations.}

\section{Education}
\cventry{2026--}{Master of Science in Machine Learning}{\href{https://uni-tuebingen.de/}{University of T\"ubingen}}{T\"ubingen, Germany}{}{Research-oriented, four-semester programme. \textbf{Master's Scholar, \href{https://eliza.school/}{Zuse School ELIZA}} (Research-Oriented Master's Scholarship). First semester, focused on the mathematical fundamentals of machine learning.}
\cventry{2025}{Exchange, Computer Science \& Engineering}{\href{https://hkust.edu.hk/}{Hong Kong University of Science and Technology}}{Hong Kong SAR}{}{Courses in machine learning, computer vision, large language models, reinforcement learning and data visualisation.}
\cventry{2023--2026}{Bachelor of Science in Engineering, Software Engineering}{\href{https://www.sdu.dk/}{University of Southern Denmark}}{S\o nderborg, Denmark}{}{Graduated June 2026. GPA \textbf{10.7/12}, thesis \textbf{12/12}. Thesis \emph{Heimdall} with Tim Lukas Adam (equal contribution), \textbf{best bachelor thesis} in a cohort of more than 100 students, chosen by the head of the programme.}

\section{Research Experience}
\cventry{Jan 2026--}{Undergraduate Student Researcher, then Research Collaborator}{\href{https://aiml-sdu.github.io}{Applied and Interpretable Machine Learning Research Group, University of Southern Denmark} (supervised by \href{https://sayvaz-mmmi.github.io/}{Serkan Ayvaz})}{S\o nderborg, Denmark; remote since Oct 2026}{}{Came up with and led my own research on what goes on inside language models, from the first question to the finished paper, designing the experiments and writing the papers with my co-authors. Still part of the group after my BSc.}
\cventry{Dec 2025--Sep 2026}{Research Collaborator}{\href{https://datavisards.com/}{DataVISards, Hong Kong University of Science and Technology} (supervised by \href{https://narechania.com/}{Arpit Narechania})}{Hong Kong SAR}{}{Helped with a study of how people fact-check data-driven news articles. It started during my exchange semester and became a paper at IEEE VIS 2026.}
\cventry{Sep 2024--Dec 2025}{Student Research Assistant}{\href{https://www.sdu.dk/en/forskning/cis/laboratories/data-and-intelligence-lab}{University of Southern Denmark Data \& Intelligence Lab}}{S\o nderborg, Denmark}{}{Built machine-learning pipelines and helped with study design, papers, grant proposals and datasets across the lab's projects.}

\section{Publications}
\cvitem{}{\small My name is in bold and $^*$ marks equal contribution. Summaries and code on \href{https://phomarkon.github.io/\#publications}{my website}.}
\subsection{Conference and Journal}
\cvitem{2026}{\textbf{P.~M.~Konrad}, T.~Tanyel, S.~Ayvaz. \emph{Routing Subspaces: Auditing Evaluation-to-Deployment Mismatch in Fine-Tuned Language Models}. \textit{Conference on Neural Information Processing Systems (NeurIPS)}, 2026, main track, poster. \iconlink{\faFile[regular]}{https://arxiv.org/abs/2607.20436}{arXiv} \enspace \iconlink{\faGlobe}{https://aiml-sdu.github.io/routing-subspaces-page/}{project} \enspace \iconlink{\faCode}{https://github.com/aiml-sdu/routing-subspaces}{code}}
\cvitem{2026}{\textbf{P.~M.~Konrad}, A.-A.~Popa, Y.~Sabzehmeidani, L.~Zhong, M.~Tripathy, A.~Constantinescu, E.~A.~Liehn, S.~Ayvaz. \emph{Challenges in Deep Learning-Based Small Organ Segmentation: A Benchmarking Perspective for Medical Research with Limited Datasets}. \textit{Biomedical Signal Processing and Control} 128, 111226, 2026. \iconlink{\faBook}{https://doi.org/10.1016/j.bspc.2026.111226}{DOI} \enspace \iconlink{\faFile[regular]}{https://arxiv.org/abs/2509.05892}{arXiv} \enspace \iconlink{\faGlobe}{https://aiml-sdu.github.io/small-organ-segmentation-page/}{project}}
\cvitem{2026}{N.-T.~Thinh$^*$, Y.~Du$^*$, \textbf{P.~M.~Konrad}, A.~Narechania. \emph{Fact-Check Your Information (FYI): A Design Probe to Understand How People Actually Fact-Check Data-Driven Articles}. \textit{IEEE VIS 2026}, \textit{IEEE Transactions on Visualization and Computer Graphics}, 2027, in press. \iconlink{\faFile[regular]}{https://arxiv.org/abs/2608.06804}{arXiv} \enspace \iconlink{\faGlobe}{https://datavisards.com/projects/2027-paper-fyi/}{project} \enspace \iconlink{\faCode}{https://github.com/DataVisards/FYI}{code}}
\cvitem{2025}{\textbf{P.~M.~Konrad}, T.~Tanyel, S.~Ayvaz. \emph{Beyond Major Floods: Deep Learning for Detecting Shallow Water Inundation in Agricultural Areas}. \textit{KES 2025}, Procedia Computer Science 270, 301--310. \iconlink{\faBook}{https://doi.org/10.1016/j.procs.2025.09.149}{DOI} \enspace \iconlink{\faGlobe}{https://aiml-sdu.github.io/beyond-major-floods-page/}{project}}
\needspace{3\baselineskip}\subsection{Workshop Papers (non-archival)}
\cvitem{2026}{\textbf{P.~M.~Konrad}, T.~Tanyel, S.~Ayvaz. \emph{Self-Reports Do Not Identify Self-Models: An Identifiability Test for Counterfactual Reports}. \textit{Philosophy of Machine Learning Workshop at ICML 2026}, oral. \iconlink{\faComments[regular]}{https://openreview.net/forum?id=iBq4yYwYfd}{OpenReview} \enspace \iconlink{\faFile[regular]}{https://arxiv.org/abs/2609.32449}{arXiv} \enspace \iconlink{\faGlobe}{https://aiml-sdu.github.io/self-reports-page/}{project}}
\cvitem{2026}{\textbf{P.~M.~Konrad}, T.~Tanyel, S.~Ayvaz. \emph{Decoded but Unused: Instruction Tuning Routes Moral Framing into the Judgment Readout}. \textit{Workshop on Mechanistic Interpretability at ICML 2026}, poster. \iconlink{\faComments[regular]}{https://openreview.net/forum?id=da9ylT2doy}{OpenReview} \enspace \iconlink{\faGlobe}{https://aiml-sdu.github.io/decoded-but-unused-page/}{project} \enspace \iconlink{\faCode}{https://github.com/aiml-sdu/decoded-but-unused}{code}}
\cvitem{2026}{\textbf{P.~M.~Konrad}, T.~Tanyel, S.~Ayvaz. \emph{A Path Already Walked: On Inheriting Network-Neuroscience Tools for Mechanistic Interpretability}. \textit{Workshop on Mechanistic Interpretability at ICML 2026}, virtual poster. \iconlink{\faComments[regular]}{https://openreview.net/forum?id=dSfjkJB3bQ}{OpenReview} \enspace \iconlink{\faGlobe}{https://aiml-sdu.github.io/path-already-walked-page/}{project} \enspace \iconlink{\faCode}{https://github.com/aiml-sdu/path-already-walked}{code}}
\needspace{3\baselineskip}\subsection{Workshop Papers (archival, in proceedings)}
\cvitem{2026}{\textbf{P.~M.~Konrad}, T.~L.~Adam, R.~Terrenzi, S.~Ayvaz. \emph{Architecture Without Architects: How AI Coding Agents Shape Software Architecture}. \textit{SAGAI Workshop at IEEE ICSA 2026}. \iconlink{\faFile[regular]}{https://arxiv.org/abs/2604.04990}{arXiv} \enspace \iconlink{\faGlobe}{https://aiml-sdu.github.io/architecture-without-architects-page/}{project} \enspace \iconlink{\faCode}{https://github.com/phomarkon/vibe-architecting-case-study}{code}}
\cvitem{2026}{T.~L.~Adam, \textbf{P.~M.~Konrad}, R.~Terrenzi, F.~G.~Lukas, R.~Yilmaz, K.~Sierszecki, S.~Ayvaz. \emph{CAKE: Cloud Architecture Knowledge Evaluation of Large Language Models}. \textit{KDA-AI Workshop at IEEE ICSA 2026}. \iconlink{\faFile[regular]}{https://arxiv.org/abs/2604.05755}{arXiv} \enspace \iconlink{\faGlobe}{https://aiml-sdu.github.io/cake/}{project} \enspace \iconlink{\faCode}{https://github.com/timadam03/CAKE-benchmark}{code}}
\cvitem{2026}{R.~Terrenzi, \textbf{P.~M.~Konrad}, T.~L.~Adam, S.~Ayvaz. \emph{A Reference Architecture for Agentic Hybrid Retrieval in Dataset Search}. \textit{SAML Workshop at IEEE ICSA 2026}. \iconlink{\faFile[regular]}{https://arxiv.org/abs/2604.16394}{arXiv}}
\needspace{3\baselineskip}\subsection{Preprints}
\cvitem{2026}{\textbf{P.~M.~Konrad}, T.~L.~Adam, A.~C.~H.~Merrild, R.~Terrenzi, R.~De~Rosa, T.~Tanyel, S.~Ayvaz. \emph{The Open-Box Fallacy: Why AI Deployment Needs a Calibrated Verification Regime}. Position paper, 2026. \iconlink{\faFile[regular]}{https://arxiv.org/abs/2605.10601}{arXiv} \enspace \iconlink{\faGlobe}{https://aiml-sdu.github.io/open-box-fallacy-page/}{project}}
\needspace{3\baselineskip}\subsection{Thesis}
\cvitem{2026}{\textbf{P.~M.~Konrad}$^*$, T.~L.~Adam$^*$. \emph{Heimdall: Only the Safe Shall Pass}. BSc thesis, Software Engineering, University of Southern Denmark. \iconlink{\faGlobe}{https://phomarkon.github.io/heimdall-page/}{project} \enspace \iconlink{\faCode}{https://github.com/phomarkon/heimdall-src}{code}}

\section{Scholarships and Awards}
\cvitem{2026}{\textbf{Research-Oriented Master's Scholarship}, \href{https://eliza.school/}{Zuse School ELIZA}.}
\cvitem{2026}{\textbf{Graduate Study Grant}, William Demant Foundation, for the MSc in T\"ubingen.}
\cvitem{2026}{\textbf{Best bachelor thesis} of 100+ students, chosen by the head of the programme, SDU (with T.~L.~Adam).}
\cvitem{2026}{\textbf{Best Startup Award finalist} (one of five), DreamBear, SDU Startup Night.}
\cvitem{2025}{\textbf{Travel grant}, Danish Data Science Academy (DDSA), to present my first paper at KES 2025.}
\cvitem{2024, 2025}{\textbf{Top 10}, Danish National Championship in AI, twice (team, then individual).}
\cvitem{2023}{\textbf{Winner}, SDU Case Competition, 30+ international teams of six from Germany and Scandinavia.}

\section{Teaching, Talks and Service}
\cvitem{Jan--Aug 2026}{\textbf{Teaching Assistant, Artificial Intelligence}, University of Southern Denmark. Led exercise sessions, helped design lab assignments and built \href{https://aiml-sdu.github.io/ai101/}{interactive course material}.}

\cvitem{Feb 2026}{{\footnotesize\faMicrophone}\,\textbf{Agentic Engineering in Practice}, invited talk, first House of Software S\o nderborg event, 100+ attendees.}

\cvitem{2025--2026}{\textbf{Student Member}, Educational Committee, Software Engineering programme, SDU.}

\section{Technical Skills}
\cvitem{Research}{Controlled experimental design, activation and path patching, linear probing, ablations, steering, fine-tuning, bootstrap uncertainty, evaluation under distribution shift, conformal prediction.}
\cvitem{ML tools}{Python, PyTorch, Hugging Face Transformers, TransformerLens, scikit-learn, pandas, NumPy, Optuna, W\&B.}
\cvitem{Engineering}{Git, Docker, Linux, SQL, Google Cloud Platform, JavaScript.}
\cvitem{Languages}{German (native), English (fluent), Danish and Thai (basic).}

\section{Other Experience}
\cvitem{2022--2026}{\textbf{Founder or co-founder} of SaturoLabs (DreamBear, an iOS bedtime-story app), Tutora ApS (CTO, tutoring platform) and Yeager GmbH (AI-assisted training app).}
\cvitem{2017--2021}{\textbf{Staff Duty Soldier}, Bundeswehr (German Armed Forces), Gl\"ucksburg. Led a small HR team for 600+ soldiers. Formal recognition and performance bonus, 2021.}

\end{document}
